\documentclass[10pt,a4paper]{amsart}
\usepackage[margin=19mm]{geometry}
\usepackage{amsmath,amssymb,mathtools}
\usepackage[scr=rsfs,cal=euler]{mathalfa}
\usepackage{xfrac}
\usepackage[unicode,hidelinks,bookmarks=false]{hyperref}
\newcommand{\ie}{i.e.\ }
\newcommand{\cf}{cf.\ }
\newcommand{\resp}{respectively\ }
\newcommand{\Subsection}{\subsection}
\providecommand{\nobreakdash}{\nobreak\hbox{-}}
\setlength{\emergencystretch}{2em}
\multlinegap0pt
\usepackage{xcolor}
\usepackage{amssymb}
\usepackage{tikz}
\usepackage[shortlabels]{enumitem}
\newtheorem{theorem}{Theorem}
\newtheorem{condition}{Condition}
\usepackage{cleveref}
\crefname{theorem}{Theorem}{Theorems}
\crefname{condition}{Condition}{Conditions}
\newcommand{\abs}[1]{\left\vert#1\right\vert}
\newcommand{\e}{\mathrm{e}}
\renewcommand{\epsilon}{\varepsilon}
\newcommand{\norm}[1]{\left\Vert#1\right\Vert}
\newcommand{\set}[1]{\left\{#1\right\}}
\newcommand{\tp}{\tuple}
\newcommand{\tuple}[1]{\left(#1\right)}
\title{Rapid Mixing for Multi-Spin Systems \\ on Girth-Five Graphs}
\author{Xiaoyu Chen, Kuikui Liu}
\date{Source: arXiv:2608.25491v1 (CC BY 4.0)}
\begin{document}
\maketitle

\noindent\emph{Source and licence.} Rapid Mixing for Multi-Spin Systems \\ on Girth-Five Graphs. Version arXiv:2608.25491v1 (CC BY 4.0), distributed by arXiv under the Creative Commons Attribution 4.0 International licence, http://creativecommons.org/licenses/by/4.0/, as recorded in the arXiv licence metadata for this exact version and shown on the abstract page https://arxiv.org/abs/2608.25491v1. Full licence notice: \texttt{licenses/arxiv-2608.25491-CC-BY-4.0.txt}.\par\smallskip

\noindent\textit{Editorial scope.} Counted unit: the article's theorem thm:Potts (source0.tex lines 329--361). The article's complete original proof of it is delivered with it (source0.tex lines 2729--2946, one \texttt{proof} environment of 218 lines, which the article places away from the statement). No mathematics is written, repaired or re-derived: the statement and the proof are the article's own lines, in the article's own notation, and no retained line is substituted or re-flowed.  Retained uncounted as the source-local context the proof needs: lines 379--398; lines 400--403; lines 848--854.  Nothing was omitted from any retained span, so the package carries no omission record.  No retained line contains a citation, so the wrapper carries no bibliography and no bibliography entry.  No retained line contains a figure environment, an image-inclusion command or drawing code, so no image is delivered and the package declares no asset dependency.  Editorial formatting only, in addition: an AMS \texttt{amsart} preamble that restates the article's own declarations and macros used by the retained lines, namely the environments \texttt{theorem}, \texttt{condition} on separate counters, together with the standard packages those lines need.  The retained environments print in this document as Theorem 1, Condition 1 in order of appearance, whereas the article prints the same items with the numbering of its own full document.  The complete native source of the article as distributed in the arXiv e-print for this exact version is bundled unchanged as \texttt{sources/208582/source0.tex}.
\begin{condition}\label{cond-local-spectral-contraction}
  Let $\epsilon \in (0,1]$ and $\alpha \geq 1$.
  We say $\mu_G$ has \emph{local $(\alpha,\epsilon)$-spectral contraction} if, with $B := \*1\*1^\intercal - A$, the following conditions hold for every $v \in V$ and every feasible pinning $\tau \in [q]^{V\setminus N[v]}$:
  \begin{itemize}
  \item it has \emph{$\alpha$-bounded marginals}:
    \begin{align*}
      \norm{B\mu^\tau_{G,v}}_{\infty} \leq \alpha/\Delta,
      \quad \text{and} \quad \norm{B\mu^\tau_{G-v,i}}_\infty \leq \alpha/\Delta, \quad \forall i \in N(v);
    \end{align*}
  \item it has \emph{$\epsilon$-contraction}:
    \begin{align*}
        \norm{
        \-{diag}\tp{\sum_{i\in N(v)} \mu^\tau_{G-v,i}}^{1/2}
        B
        \,\-{diag}\tp{\mu^\tau_{G,v}}^{1/2}
        }_{2}^2 \leq \frac{1-\epsilon}{\Delta},
    \end{align*}
  \end{itemize}
  where we treat marginal probabilities $\mu^\tau_{G,v}$ and $\mu^\tau_{G-v,i}$ as vectors in $\mathbb{R}^q$ and all norms above are  defined in Euclidean space.
\end{condition}

\begin{theorem}\label{thm:general}
  Suppose $G$ is an $n$-vertex simple graph with maximum degree $\Delta$ and girth at least $5$.
  If $\mu_G$ has local $(\alpha,\epsilon)$-spectral contraction and $\Delta = \Omega_{\alpha,\epsilon}(1) \geq 2\alpha$ and the Glauber dynamics on $\mu_G$ is irreducible, then the Glauber dynamics on $\mu_G$ has spectral gap $\Omega_{\alpha,\epsilon}(1/n)$.
\end{theorem}

\begin{align}\label{eq:spgap-implies-mixing-time}
  t_{\mathrm{mix}}(\varepsilon)
  \leq\left\lceil\frac1\gamma\left(
    \log\frac1{2\varepsilon}
    +\frac12\log\frac1{\min_{x\in\Omega}\mu(x)}
  \right)\right\rceil.
\end{align}

\begin{theorem}%[Girth-five anti-ferromagnetic Potts]
\label{thm:Potts}
Fix $\delta\in(0,1)$.  Let $G=(V,E)$ be a finite simple graph on $n$
vertices with maximum degree $\Delta$ and girth at least $5$, let
$\beta\in[0,1]$, and let $q$ be an integer.  If $\Delta$ is
sufficiently large in terms of $\delta$ and
\begin{align*}
  q\geq(1+\delta)(1-\beta)\Delta,
\end{align*}
then the Glauber dynamics for the anti-ferromagnetic $q$-state Potts
model with interaction parameter $\beta$ is irreducible, has spectral
gap $\Omega_\delta(1/n)$, and satisfies
\begin{align} \label{eq:spectral-gap-mixing-bound}
  t_{\mathrm{mix}}(\varepsilon)
  =
  O_\delta\left(
    n^2\log q+n\log\frac1\varepsilon
  \right),
  \qquad
  \varepsilon\in(0,1),
\end{align}
when $\beta=0$.  If $0<\beta\leq1$, then
\begin{align*}
  t_{\mathrm{mix}}(\varepsilon)
  =
  O_\delta\left(
    n^2\log q+n^2\Delta\log\frac1\beta
    +n\log\frac1\varepsilon
  \right),
  \qquad
  \varepsilon\in(0,1).
\end{align*}
\end{theorem}

\begin{proof}[Proof of \Cref{thm:Potts}]
We finish the proof by applying \Cref{thm:general}.
To prove the spectral-gap bound, it suffices to verify
\Cref{cond-local-spectral-contraction}.
Put $\vartheta:=1-\beta$.  The interaction matrix and its deficit matrix
are
\begin{align*}
  A=\*1\*1^\intercal-\vartheta\mathbb I,
  \qquad
  B:=\*1\*1^\intercal-A=\vartheta\mathbb I.
\end{align*}
We verify \Cref{cond-local-spectral-contraction} with
\begin{align*}
  \alpha:=\frac1\delta,
  \qquad
  \epsilon:=\frac{\delta}{1+\delta}.
\end{align*}

Fix $v\in V$ and a feasible pinning $\tau$ outside $N[v]$, and put
$d:=\deg_G(v)$.  For $i\in N(v)$ and $\!a\in[q]$, let
$k_{i,\!a}$ be the number of pinned neighbors of $i$ outside $N[v]$
whose color is $\!a$.  Since $G$ has girth at least $5$, $N[v]$
induces a star.  Thus, writing
\begin{align*}
  p_i:=\mu^\tau_{G-v,i},
\end{align*}
we have
\begin{align*}
  p_i(\!a)
  =
  \frac{\beta^{k_{i,\!a}}}
       {\sum_{\!b=1}^q\beta^{k_{i,\!b}}}.
\end{align*}
When $\beta=0$, we use the convention $0^0=1$; feasibility of $\tau$
ensures that the denominator is positive.  Bernoulli's inequality and
$\sum_{\!b}k_{i,\!b}\leq\Delta-1$ give
\begin{align*}
  \sum_{\!b=1}^q\beta^{k_{i,\!b}}
  &=
  \sum_{\!b=1}^q(1-\vartheta)^{k_{i,\!b}}
  \geq
  q-\vartheta\sum_{\!b=1}^q k_{i,\!b}
  \geq
  q-\vartheta(\Delta-1).
\end{align*}
Consequently, if $\vartheta>0$, then the assumption on $q$ yields
\begin{align}
  \vartheta p_i(\!a)
  \leq
  \frac{\vartheta}{q-\vartheta(\Delta-1)}
  \leq
  \frac1{\delta\Delta+1}
  \leq
  \frac1{\delta\Delta}.
  \label{eq:Potts-leaf-interaction-bound}
\end{align}
The same conclusion is immediate when $\vartheta=0$.

For each $\!a\in[q]$, set
\begin{align*}
  w_{\!a}
  &:=%
  \prod_{i\in N(v)}\tp{1-\vartheta p_i(\!a)},
  &
  Z&:=\sum_{\!b=1}^q w_{\!b}.
\end{align*}
The tree recursion on the conditioned star gives
\begin{align*}
  \nu_v(\!a)
  :=
  \mu^\tau_{G,v}(\!a)
  =
  \frac{w_{\!a}}Z.
\end{align*}
Since $\prod_i(1-u_i)\geq1-\sum_i u_i$ for $u_i\in[0,1]$,
\begin{align*}
  Z
  \geq
  \sum_{\!a=1}^q\tp{1-\vartheta\sum_{i\in N(v)}p_i(\!a)}
  =q-\vartheta d.
\end{align*}
It follows, again trivially when $\vartheta=0$, that
\begin{align}
  \vartheta\nu_v(\!a)
  \leq
  \frac{\vartheta}{q-\vartheta d}
  \leq
  \frac1{\delta\Delta}.
  \label{eq:Potts-center-interaction-bound}
\end{align}
Because $B=\vartheta\mathbb I$, \eqref{eq:Potts-leaf-interaction-bound}
and \eqref{eq:Potts-center-interaction-bound} prove the
$\alpha$-bounded marginal condition.

It remains to verify the contraction condition.  The case
$\vartheta=0$ is immediate, so assume $\vartheta>0$.  For each
$\!a\in[q]$, put
\begin{align*}
  T_{\!a}:=\vartheta\sum_{i\in N(v)}p_i(\!a).
\end{align*}
By \eqref{eq:Potts-leaf-interaction-bound} and
$\log(1-u)\geq-u/(1-u)$ for $u\in[0,1)$,
\begin{align*}
  \log w_{\!a}
  &\geq
  -\tp{1+\frac1{\delta\Delta}}T_{\!a}.
\end{align*}
Since $\sum_{\!a}T_{\!a}=\vartheta d$, the arithmetic--geometric
mean inequality gives
\begin{align}
  Z
  \geq
  q\tp{\prod_{\!a=1}^q w_{\!a}}^{1/q}
  \geq
  q\exp\tp{
    -\tp{1+\frac1{\delta\Delta}}
    \frac{\vartheta d}{q}
  }.
  \label{eq:Potts-normalizer-lower-bound}
\end{align}
On the other hand, $w_{\!a}\leq\exp(-T_{\!a})$.  Therefore
$T_{\!a}\exp(-T_{\!a})\leq \e^{-1}$ and
\eqref{eq:Potts-normalizer-lower-bound} imply
\begin{align*}
  &\Delta\vartheta^2\nu_v(\!a)
    \sum_{i\in N(v)}p_i(\!a)
  \\
  &\qquad=
  \Delta\vartheta T_{\!a}\frac{w_{\!a}}Z
  \\
  &\qquad\leq
  \frac{\vartheta\Delta}{q}
  \exp\tp{
    \tp{1+\frac1{\delta\Delta}}
    \frac{\vartheta d}{q}-1
  }.
\end{align*}
For sufficiently large $\Delta$ in terms of $\delta$, we may assume
$\Delta\geq\delta^{-2}$.  Using $d\leq\Delta$ and
$\vartheta\Delta/q\leq1/(1+\delta)$, we then have
\begin{align*}
  \tp{1+\frac1{\delta\Delta}}
  \frac{\vartheta d}{q}
  \leq1.
\end{align*}
Hence, uniformly in $\!a$,
\begin{align*}
  \vartheta^2\nu_v(\!a)
  \sum_{i\in N(v)}p_i(\!a)
  \leq
  \frac1{(1+\delta)\Delta}
  =
  \frac{1-\epsilon}{\Delta}.
\end{align*}
Since $B=\vartheta\mathbb I$, the matrix in the contraction condition
is diagonal and therefore
\begin{align*}
  &\norm{
    \-{diag}\tp{\sum_{i\in N(v)}p_i}^{1/2}
    B
    \-{diag}(\nu_v)^{1/2}
  }_2^2
  \\
  &\qquad=
  \vartheta^2\max_{\!a\in[q]}\set{
    \nu_v(\!a)\sum_{i\in N(v)}p_i(\!a)
  }
  \leq
  \frac{1-\epsilon}{\Delta}.
\end{align*}
The estimates above are uniform over $v$ and all feasible pinnings
$\tau$.  Thus the Potts measure has local
$(\delta^{-1},\delta/(1+\delta))$-spectral contraction.

The Glauber dynamics is irreducible.  Indeed, when $\beta>0$ the Gibbs
measure has full support.  When $\beta=0$, the assumption on $q$ gives
$q\geq\Delta+2$ for sufficiently large $\Delta$ in terms of $\delta$,
and the proper-coloring Glauber dynamics is irreducible in this range.
After increasing the lower bound on $\Delta$ once more if necessary,
all the hypotheses of \Cref{thm:general} hold.  We conclude that the
discrete-time Glauber dynamics has spectral gap
$\Omega_\delta(1/n)$.

Finally, consider the smallest positive mass of a configuration.  If
$\beta=0$, the measure is uniform on a set of at most $q^n$ proper
colorings, so
\begin{align*}
  \min_{\sigma\in\-{supp}(\mu)}\mu(\sigma)\geq q^{-n}.
\end{align*}
If $0<\beta\leq1$, then every configuration has unnormalized weight at
least $\beta^{\abs{E}}\geq\beta^{n\Delta/2}$, while the partition
function is at most $q^n$.  Hence
\begin{align*}
  \min_{\sigma\in[q]^V}\mu(\sigma)
  \geq
  \beta^{n\Delta/2}q^{-n}.
\end{align*}
Substituting these bounds and the spectral-gap estimate into the
standard worst-start mixing bound
\eqref{eq:spgap-implies-mixing-time} gives, respectively,
\begin{align*}
  t_{\mathrm{mix}}(\varepsilon)
  &=
  O_\delta\tp{
    n^2\log q+n\log\frac1\varepsilon
  },
  &&\beta=0,
  \\
  t_{\mathrm{mix}}(\varepsilon)
  &=
  O_\delta\tp{
    n^2\log q+n^2\Delta\log\frac1\beta
    +n\log\frac1\varepsilon
  },
  &&0<\beta\leq1.
\end{align*}
This proves the theorem.
\end{proof}

\end{document}
